%=============================================================
% Capitolo 2 — Sistemi Lineari
%=============================================================
\chapter{Sistemi Lineari}
\label{ch:sistemi}

\section{Introduzione}

Molti problemi applicativi in data science, ottimizzazione e simulazione si riducono
alla risoluzione di un \textbf{sistema lineare}:
\begin{equation}
  Ax = b,
  \label{eq:sistema}
\end{equation}
dove $A \in \R^{n \times n}$ è la matrice dei coefficienti, $b \in \R^n$ è il termine
noto e $x \in \R^n$ è il vettore incognito da determinare.

%------------------------------------------------------------
\section{Richiami di algebra lineare}
%------------------------------------------------------------

\subsection{Matrici e operazioni fondamentali}

Una matrice $A \in \R^{m \times n}$ ha $m$ righe e $n$ colonne. L'elemento in
posizione $(i,j)$ si indica con $a_{ij}$.

Forma esplicita:
\begin{equation}
  A =
  \begin{pmatrix}
    a_{11} & a_{12} & \cdots & a_{1n} \\
    a_{21} & a_{22} & \cdots & a_{2n} \\
    \vdots & \vdots & \ddots & \vdots \\
    a_{m1} & a_{m2} & \cdots & a_{mn}
  \end{pmatrix}.
\end{equation}

La \textbf{matrice trasposta} $A^\top \in \R^{n \times m}$ ha elementi $(A^\top)_{ij} = a_{ji}$.
La \textbf{matrice coniugata trasposta} (o aggiunta hermitiana) è $A^* = \bar{A}^\top$.

\subsection{Range e rango}

\begin{definition}[Range e nucleo]
Sia $A \in \R^{m \times n}$.
\begin{itemize}
  \item Il \emph{range} (o immagine) di $A$ è $\mathcal{R}(A) = \{Ax : x \in \R^n\} \subseteq \R^m$.
  \item Il \emph{nucleo} di $A$ è $\mathcal{N}(A) = \{x \in \R^n : Ax = 0\} \subseteq \R^n$.
\end{itemize}
\end{definition}

\begin{definition}[Rango]
Il \emph{rango} di $A$, indicato $\operatorname{rank}(A)$, è la dimensione di
$\mathcal{R}(A)$.
\end{definition}

\subsection{Determinante}

Per una matrice quadrata $A \in \R^{n \times n}$:
\begin{itemize}
  \item $\det(A) \neq 0 \iff A$ è non singolare (invertibile);
  \item Se $A$ è triangolare, $\det(A) = \prod_{i=1}^n a_{ii}$;
  \item $\det(AB) = \det(A)\det(B)$.
\end{itemize}

\subsection{Matrice inversa}

Se $A$ è non singolare, esiste unica $A^{-1}$ tale che $AA^{-1} = A^{-1}A = I$.

\begin{remark}
Calcolare $A^{-1}$ esplicitamente è \emph{costoso} ($\mathcal{O}(n^3)$) e spesso
\emph{mal condizionato}: in pratica si risolve $Ax = b$ senza mai calcolare $A^{-1}$.
\end{remark}

\subsection{Autovalori e autovettori}

\begin{definition}
$\lambda \in \mathbb{C}$ è un \emph{autovalore} di $A$ e $v \neq 0$ il corrispondente
\emph{autovettore} se
\begin{equation}
  Av = \lambda v.
  \label{eq:autovalore}
\end{equation}
\end{definition}

Gli autovalori sono le radici del \emph{polinomio caratteristico} $\det(A - \lambda I) = 0$.

\begin{definition}[Raggio spettrale]
Il \emph{raggio spettrale} di $A$ è:
\begin{equation}
  \rho(A) = \max_i |\lambda_i|,
  \label{eq:raggio_spettrale}
\end{equation}
dove $\lambda_1, \ldots, \lambda_n$ sono gli autovalori di $A$.
\end{definition}

%------------------------------------------------------------
\section{Norme vettoriali}
\label{sec:norme_vettoriali}
%------------------------------------------------------------

\begin{definition}[Norma vettoriale]
Una funzione $\|\cdot\| : \R^n \to \R$ è una \emph{norma} se:
\begin{enumerate}
  \item $\|x\| \geq 0$ e $\|x\| = 0 \iff x = 0$ \quad (positività);
  \item $\|\alpha x\| = |\alpha|\,\|x\|$ per ogni $\alpha \in \R$ \quad (omogeneità);
  \item $\|x + y\| \leq \|x\| + \|y\|$ \quad (disuguaglianza triangolare).
\end{enumerate}
\end{definition}

Le norme $p$ più comuni ($p \geq 1$) sono:
\begin{align}
  \|x\|_1   &= \sum_{i=1}^n |x_i|, \label{eq:norma1}\\
  \|x\|_2   &= \sqrt{\sum_{i=1}^n x_i^2}, \label{eq:norma2}\\
  \|x\|_\infty &= \max_{1 \leq i \leq n} |x_i|. \label{eq:normaInf}
\end{align}

\begin{theorem}[Equivalenza delle norme in $\R^n$]
Tutte le norme su $\R^n$ sono equivalenti: per ogni coppia di norme
$\|\cdot\|_\alpha$, $\|\cdot\|_\beta$ esistono costanti $c_1, c_2 > 0$ tali che
\begin{equation}
  c_1\,\|x\|_\alpha \leq \|x\|_\beta \leq c_2\,\|x\|_\alpha \quad \forall x \in \R^n.
\end{equation}
\end{theorem}

\subsection{Disuguaglianza di Cauchy-Schwarz}

\begin{theorem}[Cauchy-Schwarz]
Per ogni $x, y \in \R^n$:
\begin{equation}
  |x^\top y| \leq \|x\|_2\,\|y\|_2.
  \label{eq:cauchy_schwarz}
\end{equation}
\end{theorem}

%------------------------------------------------------------
\section{Norme matriciali}
\label{sec:norme_matriciali}
%------------------------------------------------------------

\begin{definition}[Norma matriciale indotta]
La norma matriciale \emph{indotta} (o subordinata) dalla norma vettoriale $\|\cdot\|$
è:
\begin{equation}
  \|A\| = \sup_{x \neq 0} \frac{\|Ax\|}{\|x\|} = \max_{\|x\|=1} \|Ax\|.
  \label{eq:norma_mat}
\end{equation}
\end{definition}

Proprietà fondamentale: $\|Ax\| \leq \|A\|\,\|x\|$ (compatibilità).

Le norme indotte più usate per $A \in \R^{m \times n}$:
\begin{align}
  \|A\|_1      &= \max_{1 \leq j \leq n} \sum_{i=1}^m |a_{ij}|
                  \quad \text{(massima somma di colonna)}, \\
  \|A\|_\infty &= \max_{1 \leq i \leq m} \sum_{j=1}^n |a_{ij}|
                  \quad \text{(massima somma di riga)}, \\
  \|A\|_2      &= \sigma_{\max}(A)
                  \quad \text{(massimo valore singolare)}.
\end{align}

\begin{proposition}
Per ogni norma matriciale indotta vale:
\begin{equation}
  \rho(A) \leq \|A\|.
\end{equation}
\end{proposition}

%------------------------------------------------------------
\section{Numero di condizionamento}
\label{sec:condizionamento_mat}
%------------------------------------------------------------

\begin{definition}[Numero di condizionamento]
Per una matrice non singolare $A$, il \emph{numero di condizionamento} rispetto alla
norma $\|\cdot\|$ è:
\begin{equation}
  \kappa(A) = \|A\|\,\|A^{-1}\|.
  \label{eq:cond}
\end{equation}
\end{definition}

\begin{remark}
Poiché $1 = \|I\| = \|AA^{-1}\| \leq \|A\|\,\|A^{-1}\|$, si ha sempre
$\kappa(A) \geq 1$.
\begin{itemize}
  \item $\kappa(A) \approx 1$: matrice \emph{ben condizionata}.
  \item $\kappa(A) \gg 1$: matrice \emph{mal condizionata}.
\end{itemize}
\end{remark}

Il significato pratico di $\kappa(A)$ verrà approfondito nel
\cref{ch:condizionamento}: in sintesi, $\kappa(A)$ indica di quante cifre
significative si può perdere nella soluzione di $Ax = b$ a causa degli errori di
arrotondamento.
